\clearpage
\phantomsection
\addcontentsline{toc}{section}{Nomenclatura}
\section*{Nomenclatura}

\begin{description}[leftmargin=3.2em,style=nextline]
  \item[$\mathcal R$] Conjunto de robots móviles autónomos, con $N=|\mathcal R|$.
  \item[$\mathcal L$] Conjunto de cargas o tareas físicas, con $M=|\mathcal L|$.
  \item[$q_i,\theta_i$] Posición y orientación del robot $i$.
  \item[$q_k^L$] Pose de referencia de la carga $k$.
  \item[$C_k$] Coalición de robots asignada a la carga $k$.
  \item[$n_k$] Cardinalidad mínima requerida por una carga en el modelo escalar.
  \item[$D_k$] Demanda física de la carga $k$.
  \item[$S_k$] Capacidad efectiva agregada entregada por una coalición.
  \item[$G_C$] Matriz de contactos que transforma esfuerzos locales en fuerza y torque planar.
  \item[$w_k^{\rm dem}$] \emph{Wrench} demandado por la carga: fuerza longitudinal, fuerza lateral y torque.
  \item[$\lambda$] Vector de esfuerzos admisibles en los contactos.
  \item[$\delta_k$] Déficit físico-económico de una carga, definido como demanda no satisfecha.
  \item[$y_{ik}$] Preferencia continua del robot $i$ por la carga $k$ dentro de una dinámica poblacional.
  \item[$p_{ik}$] Payoff local percibido por el robot $i$ para atender la carga $k$.
  \item[$G(t)$] Grafo de comunicación local entre robots.
  \item[AMR] \emph{Autonomous Mobile Robot}.
  \item[CBBA] \emph{Consensus-Based Bundle Algorithm}.
  \item[CBF] \emph{Control Barrier Function}.
  \item[CTDE] \emph{Centralized Training with Decentralized Execution}.
  \item[MARL] \emph{Multi-Agent Reinforcement Learning}.
  \item[MRTA] \emph{Multi-Robot Task Allocation}.
  \item[SP] Subproblema experimental de la escalera de validación.
\end{description}